\section*{Data availability}
\addcontentsline{toc}{section}{Data availability}

All data behind the paper's primary results are \emph{generated} by the
quasi-static OpenDSS simulation with synthetic event and measurement
overlays described in Section~\ref{sec:testbed}, released as part of the
software artifact together with the unmodified public test-feeder models
it consumes \cite{kersting2001} (BSD-3-Clause licence). One third-party
dataset is used, and only for the real-trace sanity check of
Section~\ref{sec:res-real}: the public MSU/ORNL power-system attack and
disturbance dataset, used unmodified under its published terms; it
supports no model or defense comparison in this paper. The federations
behind every reported synthetic result are reproducible from source with
a fixed seed:

\begin{quote}\scriptsize\ttfamily\raggedright
python -m experiments.stage\_data --out data\_battery\\
python -m experiments.run\_all --dry-run\\
python -m experiments.gen\_temporal\_battery\\
python -m experiments.gen\_magnitude\_sweep\\
python -m experiments.manifest\_tau\\
python -m experiments.gen\_gap\_x3\\
python -m experiments.run\_all --data-dir data\_battery\\
\mbox{}\quad --out-dir runs --device cuda\\
python -m experiments.report --runs-dir runs\_snn\_fixed runs\\
\mbox{}\quad runs\_x3 runs\_ssm\_x3 runs\_gap\_x3\\
\mbox{}\quad --tables-dir paper/tables --summary runs/summary.json
\end{quote}

The first regenerates the labelled per-substation windows; the second
enumerates every pre-registered run with its configuration hash; the next
four re-emit the post-registration batteries as committed manifests (the
chronological split, the magnitude, threshold and X3 gap sweeps; the
remaining X3 manifests ship committed like the rest), all regenerating
canonically; the last two execute the battery and rebuild the merged
corpus, the tables and the statistical battery, deduplicating first by
\texttt{run\_id} and then on a semantic configuration key. The generator,
the generated federations, all per-run result records, the run manifest
and the figure and table scripts are deposited in the public archive
identified below, so that the tables and figures of
Section~\ref{sec:results} can be rebuilt without re-running the battery.

\section*{Code availability}
\addcontentsline{toc}{section}{Code availability}

The complete framework -- simulation, models, federated engine, attack
suite, consensus pipeline, experiment runner, analysis code and tests --
is open source, together with the generated federations and every per-run
result record, at \url{https://github.com/ufxa/gridspike} (code under
MIT, data and results under CC BY 4.0); an archived release with a
persistent identifier will be cited here on acceptance. The
package keeps its development name (\texttt{fedgrid}) in the repository;
the system is named \sysname{} in this manuscript because that name is
already in use for an unrelated published system.

\section*{Declaration of competing interest}
\addcontentsline{toc}{section}{Declaration of competing interest}

\TODOnum{competing-interest statement: authors to confirm on the title page}
The authors declare that they have no known competing financial interests or
personal relationships that could have appeared to influence the work reported
in this paper. (The definitive statement, with author names, is on the separate
title page.)

\section*{Funding}
\addcontentsline{toc}{section}{Funding}

\TODOnum{funding sources: authors to confirm on the title page}
The definitive funding statement, including grant numbers where applicable, is
on the separate title page. If the work received no specific grant from funding
agencies in the public, commercial or not-for-profit sectors, that is stated
there.

\section*{CRediT authorship contribution statement}
\addcontentsline{toc}{section}{CRediT authorship contribution statement}

Author contributions are recorded using the CRediT taxonomy on the separate,
non-anonymized title page, as required by the double-anonymized review model.

\section*{Declaration of generative AI and AI-assisted technologies in the manuscript preparation process}
\addcontentsline{toc}{section}{Declaration of generative AI and AI-assisted technologies}

During the preparation of this work, the author(s) used
\TODOnum{name of tool / service} in order to
\TODOnum{reason: e.g.\ draft and edit sections of the manuscript}. After using
this tool/service, the author(s) reviewed and edited the content as needed and
take(s) full responsibility for the content of the published article.

\medskip
\noindent\emph{Note.} Use of artificial intelligence \emph{within the research
methodology itself} is disclosed separately in
Section~\ref{sec:ai-methodology}, in accordance with the publisher's policy that
methodological AI use belongs in the Methods rather than in this statement. No
AI system is listed as an author, and no AI-generated image is presented as
primary research data.
